\begin{remark}[counterexample; the literal blow-up shape]\label{push:src:A.2.3}\label{push:A23}
The unrestricted request that every normalized component
be a blow-up of $R$ is false. The first maximal-cone
star already contributes a $\mathbb P^3$ component,
which admits no proper bimeromorphic map to $R$, since
such a map would inject a rank-eleven rational $H^2$
into a rank-one group. \ref{push:src:A.2.4} gives every normalization
by explicit star charts; \ref{push:src:A.2.5} supplies the terminating
blow-up and blow-down algorithm with centres, order,
exceptional bundles and gluing maps. Thus the literal
shape is refuted and its general algorithmic
replacement is proved.
\end{remark}
\Uses{\cref{push:src:A.2.1,push:src:A.2.4,push:src:A.2.5}; \cref{fam:src:E.1.1}}
\begin{proof}

The first maximal-cone star of \ref{push:src:A.2.1} is the
ordinary point star in a smooth four-dimensional toric chart.
Its new central component is the exceptional divisor
$\mathbb P^3$. The specified normalization $R$ is the
projective-line bundle over the rational elliptic surface $S$
of \ref{fam:src:E.1.1}. Its integral projective-bundle basis
gives $b_2(R)=b_2(S)+1=10+1=11$, while
$b_2(\mathbb P^3)=1$.

A proper bimeromorphic map $f:\mathbb P^3\to R$ would have
topological degree one. For $a\in H^2(R;\Q)$ and
$b\in H^4(R;\Q)$, the projection formula gives
$\int_{\mathbb P^3}f^*a\,f^*b=\int_Rab$.
Poincare duality makes $f^*$ injective in degree two, contradicting
these ranks. Thus this explicit component refutes the proposed
blow-up-only presentation. The star charts and the terminating
modification sequence in \ref{push:src:A.2.4,push:src:A.2.5}
give the stated replacement for every allowed fan.

\end{proof}
